%----------------------------------------------------------------------------------------
%	PACKAGES AND OTHER DOCUMENT CONFIGURATIONS
%----------------------------------------------------------------------------------------

\usepackage{amsmath,amsfonts,stmaryrd,amssymb} % Math packages

%----------------------------------------------------------------------------------------
%  COLOR PALETTE (for listings, boxes, headings)
%----------------------------------------------------------------------------------------

\usepackage{xcolor}

% Palette intentionally different from the default look
\definecolor{NCISAccent}{HTML}{1F4E79}   % deep blue
\definecolor{NCISAccent2}{HTML}{7F2B3A}  % burgundy
\definecolor{NCISInk}{HTML}{1F2328}      % near-black
\definecolor{NCISMuted}{HTML}{6A737D}    % muted gray
\definecolor{NCISBg}{HTML}{F6F8FA}       % light background
\definecolor{NCISCodeBg}{HTML}{F3F6FA}   % code background
\definecolor{NCISFrame}{HTML}{D0D7DE}    % subtle border

\usepackage[ddmmyyyy]{datetime}
\usepackage{enumerate} % Custom item numbers for enumerations

\usepackage[ruled]{algorithm2e} % Algorithms

\usepackage[framemethod=tikz]{mdframed} % Allows defining custom boxed/framed environments

\usepackage{listings} % File listings, with syntax highlighting
\usepackage[hidelinks]{hyperref}

\lstset{
  language=C,
  basicstyle=\ttfamily\small\color{NCISInk},
  numbers=left,
  numberstyle=\scriptsize\color{NCISMuted},
  stepnumber=1,
  numbersep=4pt,
  frame=single,
  rulecolor=\color{NCISFrame},
  framesep=5pt,
  xleftmargin=20pt,
  framexleftmargin=15pt,
  breaklines=true,
  showstringspaces=false,
  backgroundcolor=\color{NCISCodeBg},
  keywordstyle=\color{NCISAccent}\bfseries,     % parole chiave (if, else, int, return...)
  commentstyle=\color{NCISMuted}\itshape,       % commenti
  stringstyle=\color{NCISAccent2},              % stringhe
  identifierstyle=\color{NCISInk},              % identificatori normali
  emph={size_t, uint32_t, uint64_t, TAS, RTE},            % tipi aggiuntivi
  emphstyle=\color{NCISAccent2}\bfseries,      % stile tipi aggiuntivi
  morekeywords={uint8_t, uint16_t, uint32_t, uint64_t, int8_t, int16_t, int32_t, int64_t, bool, true, false}, % tipi e parole chiave extra
  morecomment=[l]{//},                           % commenti con //
  morecomment=[s]{/*}{*/},                       % commenti multilinea
}

\lstdefinelanguage{RISCAsm}{
  morekeywords={
    ADDQ, SUM, ANDI, MOVE,MOVEA, MOVEM, DC, DS, sra, 
    lw, sw, li, la, mv, 
    BEQ, BNE, JMP, CMP, jalr, ret, 
    NOP, lui, RTE, RTS, 
    ecall, ebreak, ORG, EQU, JSR, CLR, TAS 
  },
  sensitive=true,
  morecomment=[l]{\*},
  morestring=[b]",
  basicstyle=\ttfamily\small\color{NCISInk},
  keywordstyle=\color{NCISAccent}\bfseries,
  commentstyle=\color{NCISMuted}\itshape,
  stringstyle=\color{NCISAccent2},
  numbers=left,
  numberstyle=\scriptsize\color{NCISMuted},
  stepnumber=1,
  numbersep=4pt,
  frame=single,
  rulecolor=\color{NCISFrame},
  framesep=5pt,
  xleftmargin=20pt,
  framexleftmargin=15pt,
  showstringspaces=false,
  breaklines=true,
  backgroundcolor=\color{NCISCodeBg},
}

%----------------------------------------------------------------------------------------
% PYTHON LISTING STYLE
%----------------------------------------------------------------------------------------

\lstdefinestyle{PythonStyle}{
  language=Python,
  basicstyle=\ttfamily\footnotesize\color{NCISInk},
  numbers=left,
  numberstyle=\scriptsize\color{NCISMuted},
  stepnumber=1,
  numbersep=4pt,
  frame=single,
  rulecolor=\color{NCISFrame},
  framesep=5pt,
  xleftmargin=20pt,
  framexleftmargin=15pt,
  breaklines=true,
  showstringspaces=false,
  backgroundcolor=\color{NCISCodeBg},
  keywordstyle=\color{NCISAccent}\bfseries,
  commentstyle=\color{NCISMuted}\itshape,
  stringstyle=\color{NCISAccent2},
  % Inclusione di keyword comuni e built-in
  morekeywords={as, assert, break, class, continue, def, del, elif, else, except, False, finally, for, from, global, if, import, in, is, lambda, None, nonlocal, not, or, pass, raise, return, True, try, while, with, yield, self},
  % Stile per i numeri
  emph={1, 2, 3, 4, 5, 6, 7, 8, 9, 0},
  emphstyle=\color{NCISInk}
}

% Ambiente personalizzato per snippet Python veloci
\lstnewenvironment{pythoncode}[1][]
{
\lstset{style=PythonStyle,#1}
}{}

%----------------------------------------------------------------------------------------
%	DOCUMENT MARGINS
%----------------------------------------------------------------------------------------

\usepackage{geometry} % Required for adjusting page dimensions and margins

\geometry{
	paper=a4paper, % Paper size, change to letterpaper for US letter size
	top=2.5cm, % Top margin
	bottom=3cm, % Bottom margin
	left=2.5cm, % Left margin
	right=2.5cm, % Right margin
	headheight=14pt, % Header height
	footskip=1.5cm, % Space from the bottom margin to the baseline of the footer
	headsep=1.2cm, % Space from the top margin to the baseline of the header
	%showframe, % Uncomment to show how the type block is set on the page
}

%----------------------------------------------------------------------------------------
%	FONTS
%----------------------------------------------------------------------------------------

\usepackage[utf8]{inputenc} % Required for inputting international characters
\usepackage[T1]{fontenc} % Output font encoding for international characters

% Font: prefer a distinct, modern Times-like face (with safe fallback)
\IfFileExists{newtxtext.sty}{%
  \usepackage{newtxtext}
  \usepackage{newtxmath}
}{%
  \usepackage{XCharter}
}

\usepackage{microtype}
\usepackage{enumitem}
\usepackage{tikz} 
\usetikzlibrary{positioning}

%----------------------------------------------------------------------------------------
%  PAGE STYLE + HEADINGS (purely aesthetic)
%----------------------------------------------------------------------------------------

\usepackage{fancyhdr}
\usepackage{titlesec}
\usepackage{titling}

% Slightly more modern paragraph rhythm
\setlength{\parindent}{0pt}
\setlength{\parskip}{0.55em}

% Header/footer
\pagestyle{fancy}
\fancyhf{}
\renewcommand{\headrulewidth}{0.6pt}
\renewcommand{\headrule}{\hbox to\headwidth{\color{NCISFrame}\leaders\hrule height \headrulewidth\hfill}}
\fancyhead[L]{\color{NCISMuted}\small\nouppercase{\leftmark}}
\fancyhead[R]{\color{NCISMuted}\small\thepage}

\fancypagestyle{plain}{%
  \fancyhf{}
  \renewcommand{\headrulewidth}{0pt}
  \fancyfoot[C]{\color{NCISMuted}\small\thepage}
}

% Section title styling
\titleformat{\section}
  {\Large\bfseries\color{NCISAccent}}
  {\thesection}{0.75em}{}
  [\vspace{0.25em}\titlerule]

\titleformat{\subsection}
  {\large\bfseries\color{NCISAccent2}}
  {\thesubsection}{0.75em}{}

\titlespacing*{\section}{0pt}{1.2em}{0.6em}
\titlespacing*{\subsection}{0pt}{0.9em}{0.4em}

% Title block (maketitle)
\pretitle{\vspace{-1em}\begin{center}\LARGE\bfseries\color{NCISAccent}}
\posttitle{\par\vspace{0.4em}\end{center}}
\preauthor{\begin{center}\large\color{NCISInk}}
\postauthor{\par\end{center}}
\predate{\begin{center}\small\color{NCISMuted}}
\postdate{\par\vspace{1.2em}\end{center}}

% pseudocode environment %


%----------------------------------------------------------------------------------------
%	COMMAND LINE ENVIRONMENT
%----------------------------------------------------------------------------------------

% Usage:
% \begin{commandline}
%	\begin{verbatim}
%		$ ls
%		
%		Applications	Desktop	...
%	\end{verbatim}
% \end{commandline}

\mdfdefinestyle{commandline}{
	leftmargin=10pt,
	rightmargin=10pt,
	innerleftmargin=15pt,
  middlelinecolor=NCISFrame,
	middlelinewidth=2pt,
	frametitlerule=false,
  backgroundcolor=NCISBg,
	frametitle={Command Line},
  frametitlefont={\normalfont\sffamily\bfseries\color{white}\hspace{-0.8em}},
  frametitlebackgroundcolor=NCISAccent,
	nobreak,
}

% Define a custom environment for command-line snapshots
\newenvironment{commandline}{
	\medskip
	\begin{mdframed}[style=commandline]
}{
	\end{mdframed}
	\medskip
}

%----------------------------------------------------------------------------------------
%	FILE CONTENTS ENVIRONMENT
%----------------------------------------------------------------------------------------

% Usage:
% \begin{file}[optional filename, defaults to "File"]
%	File contents, for example, with a listings environment
% \end{file}

\mdfdefinestyle{file}{
    innertopmargin=1.6\baselineskip,
    innerbottommargin=0.8\baselineskip,
    topline=false, bottomline=false,
    leftline=false, rightline=false,
    leftmargin=2cm,
    rightmargin=2cm,
    singleextra={
        \draw[fill=black!10!white](P)++(0,-1.2em)rectangle(P-|O);
        \node[anchor=north west] at(P-|O){\ttfamily\mdfilename};
        \def\l{3em}
        \draw(O-|P)++(-\l,0)--++(\l,\l)--(P)--(P-|O)--(O)--cycle;
        \draw(O-|P)++(-\l,0)--++(0,\l)--++(\l,0);
    },
}

% Comando per tab a 4 spazi all’interno dell’ambiente file
\newenvironment{file}[1][File]{%
    \medskip
    \newcommand{\mdfilename}{#1}
    % Rende il carattere tab attivo e lo definisce come 4 spazi
    \catcode`\^^I=\active
    \def^^I{\hspace*{4ex}}%
    \begin{mdframed}[style=file]
}{%
    \end{mdframed}
    \medskip
}

%----------------------------------------------------------------------------------------
%	NUMBERED QUESTIONS ENVIRONMENT
%----------------------------------------------------------------------------------------

% Usage:
% \begin{question}[optional title]
%	Question contents
% \end{question}

\mdfdefinestyle{question}{
	innertopmargin=1.2\baselineskip,
	innerbottommargin=0.8\baselineskip,
  roundcorner=7pt,
  linecolor=NCISFrame,
  linewidth=0.8pt,
  backgroundcolor=NCISBg,
	nobreak,
	singleextra={%
		\draw(P-|O)node[xshift=1em,anchor=west,fill=white,draw,rounded corners=5pt]{%
		Question \theQuestion\questionTitle};
	},
}

\newcounter{Question} % Stores the current question number that gets iterated with each new question

% Define a custom environment for numbered questions
\newenvironment{question}[1][\unskip]{
	\bigskip
	\stepcounter{Question}
	\newcommand{\questionTitle}{~#1}
	\begin{mdframed}[style=question]
}{
	\end{mdframed}
	\medskip
}

%----------------------------------------------------------------------------------------
%	WARNING TEXT ENVIRONMENT
%----------------------------------------------------------------------------------------

% Usage:
% \begin{warn}[optional title, defaults to "Warning:"]
%	Contents
% \end{warn}

\mdfdefinestyle{warning}{
	topline=false, bottomline=false,
	leftline=false, rightline=false,
	nobreak,
	singleextra={%
    \fill[NCISAccent2](P-|O)rectangle++(0.9em,-0.9em);
    \node[anchor=center] at ($(P-|O)+(0.45em,-0.45em)$){\color{white}\scriptsize\bfseries !};
    \draw[very thick,NCISAccent2](P-|O)++(0,-1em)--(O);
	}
}

% Define a custom environment for warning text
\newenvironment{warn}[1][Warning:]{ % Set the default warning to "Warning:"
	\medskip
	\begin{mdframed}[style=warning]
		\noindent{\textbf{#1}}
}{
	\end{mdframed}
}

%----------------------------------------------------------------------------------------
%	INFORMATION ENVIRONMENT
%----------------------------------------------------------------------------------------

% Usage:
% \begin{info}[optional title, defaults to "Info:"]
% 	contents
% 	\end{info}

\mdfdefinestyle{info}{%
	topline=false, bottomline=false,
	leftline=false, rightline=false,
	nobreak,
	singleextra={%
    \fill[NCISAccent](P-|O)circle[radius=0.42em];
    \node at(P-|O){\color{white}\scriptsize\bfseries i};
    \draw[very thick,NCISAccent](P-|O)++(0,-0.8em)--(O);
	}
}

% Define a custom environment for information
\newenvironment{info}[1][Info:]{ % Set the default title to "Info:"
	\medskip
	\begin{mdframed}[style=info]
		\noindent{\textbf{#1}}
}{
	\end{mdframed}
}
